%%
%% 张哲简历
%%

\documentclass[zh]{resume}

\setmainfont{TeX Gyre Termes}
\setmonofont{Noto Sans Mono}
\setCJKmainfont{Noto Serif CJK SC}

\fileinfo{%
  \faCopyright{} 2026, Zhang Zhe \hspace{0.5em}
  \faEdit{} \today
}

\name{哲}{张}
\makeatletter
\def\@photo@width{2.35cm}
\def\@photo@position{right}
\def\@photo{%
  \begin{minipage}{\@photo@width}
    \includegraphics[width=\linewidth]{zhangzhe.jpg}%
  \end{minipage}%
}
\makeatother

\keywords{Agent 系统, LangGraph, MCP, RAG, Django, 智能运维, Vue3, Node.js}

\profile{
  \mobile{13160006317}
  \email{3220857484@qq.com}
  \university{南京理工大学}
  \degree{软件工程 \textbullet 硕士在读}
}

\begin{document}
\makeheader

\begin{abstract}
南京理工大学软件工程硕士在读，本科毕业于南通大学软件工程专业。方向聚焦 Agent 系统与 LLM 应用工程，熟悉 LangGraph/LangChain 工作流编排、MCP 工具接入、RAG 检索增强和评估链路设计，能够完成从状态图设计、工具服务开发、后端接口到前端可视化的完整落地。主导实现 OCR 运维诊断 Agent 和教辅录题审核系统，并在模型微调、检索优化方向具备竞赛实践经验。
\end{abstract}

\sectionTitle{教育背景}{\faGraduationCap}
\begin{educations}
  \education%
    {2024.09}%
    [至今]%
    {南京理工大学}%
    {软件工程专业}%
    {软件工程}%
    {硕士}
  \separator{0.5ex}
  \education%
    {2019.09}%
    [2023.06]%
    {南通大学}%
    {软件工程专业}%
    {软件工程}%
    {本科}
\end{educations}

\sectionTitle{比赛经历}{\faTrophy}
\begin{itemize}
  \item \textbf{KDD Cup 2024 学术论文问答检索挑战赛} \hfill 2024.05 -- 2024.07
  \begin{itemize}
    \item \textbf{任务背景}: 面向学术问答场景，从海量论文中检索与专业问题最匹配的候选论文。
    \item \textbf{方案设计}: 针对复杂 Query 冗长、语义匹配难的问题，设计 ``问题精炼 + 向量召回 + 重排序'' 两阶段检索方案；基于 Llama3 对 Query 进行精炼总结，使用 BGE M3 对 Query 与论文语料统一向量化，结合 FAISS 召回 Top100 候选论文，再引入 BGE-Reranker 对结果精排筛选 Top20。
    \item \textbf{优化细节}: 通过 Hard Example Mining 构建微调数据，对 Reranker 进行专项优化，提升困难样本区分能力。
    \item \textbf{结果}: 官方测试集检索分数从 0.138 提升至 0.166，较基线提升约 20.3\%。
  \end{itemize}

  \item \textbf{``天翼云息壤杯'' AI 大赛：大语言模型数学推理能力提升} \hfill 2024.11 -- 2025.07
  \begin{itemize}
    \item \textbf{任务背景}: 基于 Qwen3-8B 开展微调训练，提升模型对复杂数学问题与深度推理的理解求解能力。
    \item \textbf{方案设计}: 围绕数学推理任务长链路推导困难、监督样本质量不足的问题，设计 ``数据蒸馏 + 质量筛选 + 参数高效微调'' 整体方案；基于 MetaMathQA 构建数据集，利用 Qwen2.5-Math-72B 蒸馏生成 10 万条高质量 CoT 样本，并结合 Math-Verify 与 Qwen2.5-72B 评审模型做答案校验。
    \item \textbf{工程落地}: 在 8 卡昇腾 910B 64G 环境完成 Qwen3-8B 的 LoRA 微调，并结合 vLLM 优化推理性能，实现训推一体化部署。
    \item \textbf{结果}: 微调后的 Qwen3-8B 在客观数学题上的正确率从 38\% 提升至 62\%。
  \end{itemize}
\end{itemize}

\sectionTitle{专业技能}{\faCogs}
\begin{itemize}
  \item \textbf{Agent 开发}: 熟悉 LangGraph 状态图编排、条件路由、循环节点和状态管理，能够设计带工具调用、可追踪、可回放的 Tool-use Agent；熟悉 MCP（Model Context Protocol）协议，具备自定义工具服务开发经验。
  \item \textbf{LLM 应用}: 熟悉 LangChain 工具链与 Prompt 设计，具备 RAG 检索增强（向量召回 + 重排序）、规则 + LLM 混合决策、知识蒸馏、LoRA 微调及 GRPO 强化学习实践经验；了解 vLLM 推理部署。
  \item \textbf{后端开发}: 熟悉 Python（FastAPI / Django）和 Node.js（Fastify）后端开发，掌握 RESTful API、SSE 流式接口、PostgreSQL/MySQL、Redis、Celery/BullMQ 消息队列；具备 Django REST framework、Prisma ORM、JWT 认证经验。
  \item \textbf{前端开发}: 熟悉 Vue 3 + TypeScript + Vite 技术栈，了解 React/Next.js；掌握 Element Plus、Ant Design Vue、Pinia、动态路由等前端工程化实践，能够实现 Agent 执行过程可视化页面。
  \item \textbf{DevOps}: 熟悉 Docker/Docker Compose 容器化部署，掌握 MinIO 对象存储、Prometheus + Grafana 监控体系搭建，具备 Nginx 反向代理和基础 CI/CD 配置经验。
\end{itemize}

\sectionTitle{实习经历}{\faBriefcase}
\begin{itemize}
  \item \textbf{中国电子科技集团公司第二十八研究所（智能模型资源管理平台研发实习）} \hfill 2025.09 -- 至今
  \begin{itemize}
    \item \textbf{项目定位}: 参与智能模型管理平台开发，系统包含模型、算法、数据集管理，以及版本维护、文件交付、语义检索和部署等功能。
    \item \textbf{模型生命周期管理}: 负责模型库和版本管理模块开发，梳理模型从录入、测试、审核到上架/下架的流程，补齐分类、可见性、状态流转和历史记录等功能；将上架前检查做成规则校验，自动检查模型描述、分类、关联数据集、文档/测试报告和 Docker 镜像配置，减少人工核对。
    \item \textbf{文件与镜像交付}: 负责模型文件和镜像交付链路开发，基于 MinIO 实现预签名上传和分片上传，支持权重、说明文档、测试报告等大文件直传；基于 Celery 处理镜像导入任务，完成镜像上传、入库和版本信息同步，减少人工操作。
    \item \textbf{向量检索}: 负责模型语义检索功能开发，将模型描述、分类、关联数据集/算法和版本说明整理成检索文本，基于 Qwen3-Embedding-0.6B + FAISS 实现本地向量检索，并支持索引增量更新和删除；检索结果复用权限和可见性过滤，方便用户按任务描述快速找到合适模型。
  \end{itemize}
\end{itemize}

\sectionTitle{项目经历}{\faCode}
\begin{itemize}
  \item \textbf{OCR 运维诊断 Agent（DeerFlow / LangGraph）} \hfill 2026.01 -- 至今
  \begin{itemize}
    \item \textbf{诊断链路设计}: 负责 OCR 运维诊断 Agent 设计与落地，基于 DeerFlow / LangGraph 构建 Triage $\rightarrow$ Diagnosis $\rightarrow$ Report 三阶段工作流，支持运维人员以自然语言输入告警或排障问题，自动完成告警分类、证据收集、根因分析与处置建议生成。
    \item \textbf{工具接入与分析策略}: 接入 Prometheus 指标查询、Docker 容器状态与日志采集、历史案例检索等诊断工具，覆盖 API P95、队列 waiting/failed/workers、CPU/内存占用等关键证据；在 Diagnosis 阶段采用规则诊断 + LLM 校对策略，降低纯模型推断带来的误判风险。
    \item \textbf{可观测性与工程化}: 开发 AIOps 诊断台，基于 SSE 展示 Agent 各阶段输入输出与分析过程，支持诊断链路回放；补充接口测试与联调脚本，并完成 OCR 开发环境联调，便于故障复盘与演示。
  \end{itemize}

  \item \textbf{教辅录题审核一体化系统} \hfill 2025.06 -- 至今
  \begin{itemize}
    \item \textbf{主流程建设}: 负责录题审核主流程建设，基于 Vue 3 + TypeScript、Fastify + Prisma + PostgreSQL + Redis 开发录题工作台及后端服务，支持 PDF 双栏浏览、框选 OCR、公式渲染、草稿自动保存等能力，打通教材上传、教师录题、质检抽检、结果入库链路，替代原有 Excel + 多工具协作方式。
    \item \textbf{结构化自动切题}: 主导自动切题模块设计与落地，基于 MinerU 的块级坐标解析与目录分段能力，将教材 PDF 转换为可定位的结构化内容；结合大模型完成题号、题干、选项、答案、解析抽取，并基于 ``章节标题 + 题号'' 实现题答配对，生成包含编号、页码、坐标、题型等字段的结构化标注结果回写工作台，支撑后续校对、编辑与入库。
    \item \textbf{稳定性与工程化}: 负责长任务链路治理，基于 BullMQ 承接 PDF 拆页、MinerU 解析与自动切题任务，避免重任务阻塞录题工作台；补充队列监控告警、草稿缓存与历史恢复、提交前校验等机制，降低录题过程中的卡顿、丢稿和返工风险。
  \end{itemize}
\end{itemize}


\sectionTitle{荣誉奖项}{\faTrophy}
\begin{itemize}
  \item WACV 2025 目标距离估计挑战赛，国家级冠军
  \item ICCV 2025 感知测试挑战赛：多项选择题视频问答，国家级冠军
  \item 天翼云息壤杯 高校 AI 大模型比赛（华为）国家级三等奖
  \item CVPR 2025 少样本目标检测，国家级三等奖
\end{itemize}

\end{document}
